\section{Motivation}\label{sec:motivation}
Consider a reviewer who runs a differential check on a refactoring and sees ``no difference''. Three situations produce the same output: the changed lines ran and behaved identically; the changed lines never ran; or the comparison was unreliable (flaky output, a timeout on one side, a missing entry point). A tool that reports all three as ``pass'' invites false assurance. The ambiguity of an empty difference set was already stated at the origin of differential testing for change detection~\cite{evans2007}. Impact-analysis authors note that their results hold only for observed executions~\cite{law2003,orso2003}, and invariant-based work notes that accuracy depends on test quality~\cite{ernst2001}. Change coverage is measured in practice~\cite{hilton2018a}, but as a property of a test suite rather than as a condition on the reported outcome.

The question that motivates V7-B is narrow and empirical: if the verdict is made conditional on a recorded sufficiency level and on reproducibility checks, does the resulting ``no difference'' carry measurably less false assurance than a same-budget baseline that has no sufficiency gate? Section~\ref{sec:results} gives the exploratory answer: yes, by a margin not shown to reach the pre-specified practical threshold, and entirely through abstention.
